\section{拆掉 Language：为什么 Language 是毒药}

本章进入节目标题里的核心技术判断。很多 VLA 路线会把传感器信号先翻译成 language token，再用 language token 解码 trajectory。刘先明认为这会产生瓶颈，因为中间语言表示引入人工监督和人工语义压缩，不利于 data scaling。小鹏的做法是拆掉这个中间 language 层，让 vision 和 language 作为输入，直接解码 action。

\begin{figure}[H]
\centering
\includegraphics[width=0.94\textwidth]{remove-language-bottleneck.png}
\caption{拆掉 Language Bottleneck：中间语言 token 会限制数据 scaling。自制概念图，依据 00:19:00--00:33:53 对谈内容整理。}
\end{figure}

\begin{knowledgebox}{读图：保留语言作为输入，拆掉语言作为中间瓶颈}
左侧路线把传感器信号翻译成语言，再从语言解码轨迹；右侧路线让 Vision+Language 直接进入模型，Action 直接输出。这里不是不要语言，而是不要把语言当成强制中间监督信号。
\end{knowledgebox}

\subsection{语言为什么会变成瓶颈}

本节先把“毒药”这个强表达拆开看。语言的问题不是语义无用，而是它一旦成为低层驾驶信号的必经中介，就会把连续世界压成离散标签。

语言是人类高层抽象，非常适合指令、解释和常识，但自动驾驶的传感器信号包含大量连续、几何、动态和细粒度信息。如果必须先翻译成语言，就会丢掉许多连续信息，并把训练过程变成人工标注或人工 refinement。对需要海量数据 scaling 的自动驾驶来说，这会限制数据使用效率。

\begin{importantbox}{Language 是毒药的准确含义}
语言作为输入很有价值，语言作为强制中间表示可能有毒。毒性来自过度依赖人工语义、降低连续数据利用率、形成 bottleneck，并阻碍自监督或大规模数据训练。
\end{importantbox}

\subsection{“拆掉 L”与自监督学习}

过去 AI 成功的重要经验之一，是使用数据做 unsupervised/self-supervised learning。自监督学习利用数据自身结构构造训练信号，减少人工标注依赖。自动驾驶如果要走这条路，就要尽量减少中间人工监督层，让模型从海量视觉、传感器和驾驶轨迹中直接学习。

\begin{warningbox}{不要把“拆掉语言”理解成“不要语言能力”}
自动驾驶仍需要理解导航指令、交通语义和人类意图。拆掉的是中间瓶颈，不是人机交互里的语言输入，也不是模型中的语义理解。
\end{warningbox}

\subsection{本章小结}

“Language 是毒药”是对中间表示的警惕。物理 AI 要充分利用连续传感器数据，语言不能成为所有信息的窄口。更直接的 VLA/Action decoding，是为了提升数据使用效率和 scaling 能力。

